%! AP Statistics — Unit 4, Lesson 4.11 — Homework
\documentclass[10pt]{article}
\usepackage{apstats-article}
\usepackage{apstats-boxes}
\usepackage{grading}
\gradingassignment{apstats-u04-l11-hw}
\gradingcourse{advanced_statistics}

\begin{document}

\pageheader{Unit 4, Lesson 4.11}{Homework}

\vspace{0.1in}

\begin{practicebox}
\small
\textbf{Directions:} For each problem, show all formula setup before computing.
Write a complete-sentence interpretation wherever asked.

\vspace{10pt}
\textbf{1.} A seed company claims its tomato seeds have an 85\% germination rate.
A gardener plants 30 seeds. Let $X =$ the number that germinate.
Assume $X \sim B(30, 0.85)$.
\begin{enumerate}[label=(\alph*), leftmargin=2em, itemsep=8pt]
  \item Calculate $\mu_X$ and $\sigma_X$. Show formula setup.\\
        \writeline\\[1pt]\writeline\gradedhide{\hrulefill}\gradedright{\raisebox{-2pt}{$\mu_X =$~\gradedblank[id=hw1a1, type=short, points=1, answer=25.5, check=numeric, tolerance=0.05]{0.9in}{0.2in}\quad $\sigma_X =$~\gradedblank[id=hw1a2, type=short, points=1, answer=1.96, check=numeric, tolerance=0.01]{0.9in}{0.2in}}}
  \item Interpret $\mu_X$ in context.\\
        \gradedhide{\writeline\gradedright{\raisebox{-5pt}{\gradedblank[id=hw1b, type=rubric, points=2, display={Over many plantings of 30 seeds, an average of 25.5 seeds germinate.}]{6.1in}{\gradedlinesheight{2}-5pt}}}\\[1pt]\writeline}
  \item Find $P(X \le 22)$. Write the calculator notation.\\
        \writeline\gradedhide{\hrulefill}\gradedright{\raisebox{-2pt}{\gradedblank[id=hw1c, type=short, points=1, answer=0.0698, check=numeric, tolerance=0.002]{1.3in}{0.2in}}}
  \item Find $P(X \ge 28)$. Write the inequality rewrite and the calculator notation.\\
        \writeline\\[1pt]\writeline\gradedhide{\hrulefill}\gradedright{\raisebox{-2pt}{\gradedblank[id=hw1d, type=short, points=2, answer=0.1514, check=numeric, tolerance=0.002]{1.3in}{0.2in}}}
\end{enumerate}

\vspace{8pt}
\textbf{2.} A multiple-choice test has 50 questions, each with 5 options.
A student guesses on every question. Let $X =$ the number correct.
\begin{enumerate}[label=(\alph*), leftmargin=2em, itemsep=8pt]
  \item State the distribution of $X$: \makebox[4cm][l]{\hspace{1.6cm}\smash{\raisebox{-4pt}{\gradedblank[id=hw2a, type=short, points=1, answer={B(50, 0.20)}, check=text]{2.4cm}{0.2in}}}}
  \item Compute $\mu_X$ and $\sigma_X$.\\
        \writeline\gradedhide{\hrulefill}\gradedright{\raisebox{-2pt}{$\mu_X =$~\gradedblank[id=hw2b1, type=short, points=1, answer=10, check=numeric, tolerance=0, display={\textbf{10}}]{0.9in}{0.2in}\quad $\sigma_X =$~\gradedblank[id=hw2b2, type=short, points=1, answer=2.83, check=numeric, tolerance=0.01]{0.9in}{0.2in}}}
  \item The passing score is 30 out of 50. Find $P(X \ge 30)$.
        Is this likely to happen by guessing?\\
        \writeline\gradedhide{\hrulefill}\gradedright{\raisebox{-2pt}{\gradedblank[id=hw2c1, type=short, points=1, answer=0, check=numeric, tolerance=0.001, display={$\approx 0.0000$}]{1.3in}{0.2in}}}\\[1pt]\gradedhide{\writeline\gradedright{\raisebox{-1pt}{\gradedblank[id=hw2c2, type=rubric, points=2, display={No: 30 is about 7 SDs above the mean of 10; essentially impossible by guessing.}]{6.1in}{0.22in}}}}
\end{enumerate}

\vspace{8pt}
\textbf{3.} A hospital reports that 40\% of emergency room patients are discharged
within 4 hours. On a given day, 120 patients visit the ER.
Let $X =$ the number discharged within 4 hours. Assume $X \sim B(120, 0.40)$.
\begin{enumerate}[label=(\alph*), leftmargin=2em, itemsep=8pt]
  \item $\mu_X = $ \makebox[2cm][l]{\hspace{1.6cm}\smash{\raisebox{-4pt}{\gradedblank[id=hw3a1, type=short, points=1, answer=48, check=numeric, tolerance=0]{0.9cm}{0.2in}}}} \qquad $\sigma_X \approx$ \makebox[2cm][l]{\hspace{1.6cm}\smash{\raisebox{-4pt}{\gradedblank[id=hw3a2, type=short, points=1, answer=5.37, check=numeric, tolerance=0.01]{2cm}{0.2in}}}}
  \item Interpret $\mu_X$: \gradedhide{\writeline\gradedright{\raisebox{-5pt}{\gradedblank[id=hw3b, type=rubric, points=2, display={Over many days with 120 ER patients, an average of 48 are discharged within 4 hours.}]{4.9in}{\gradedlinesheight{2}-5pt}}}\\[1pt]\writeline}
  \item Find $P(40 \le X \le 50)$. Write this as a difference of two \texttt{binomcdf} values.\\
        \writeline\\[1pt]\writeline\gradedhide{\hrulefill}\gradedright{\raisebox{-2pt}{\gradedblank[id=hw3c, type=short, points=2, answer=0.6257, check=numeric, tolerance=0.002]{1.3in}{0.2in}}}
\end{enumerate}

\vspace{8pt}
\textbf{4.} (AP-Style Free Response) \textit{According to a survey, 62\% of adults in
a large city support a new public transit plan. A researcher randomly samples
45 adults (less than 10\% of the city's adult population).
Let $X =$ the number who support the plan.}
\begin{enumerate}[label=(\alph*), leftmargin=2em, itemsep=8pt]
  \item Define the distribution of $X$ and check all BINS conditions briefly.\\
        \writeline\\[1pt]\gradedhide{\writeline\gradedright{\raisebox{-5pt}{\gradedblank[id=hw4a, type=rubric, points=3, display={$X \sim B(45, 0.62)$: support/not; sample $< 10\%$ of adults so independent; $n = 45$ fixed; $p = 0.62$ each.}]{6.1in}{\gradedlinesheight{2}-5pt}}}\\[1pt]\writeline}
  \item Calculate and interpret $\mu_X$.\\
        \writeline\gradedhide{\hrulefill}\gradedright{\raisebox{-2pt}{\gradedblank[id=hw4b1, type=short, points=1, answer=27.9, check=numeric, tolerance=0.05]{1.3in}{0.2in}}}\\[1pt]\gradedhide{\writeline\gradedright{\raisebox{-1pt}{\gradedblank[id=hw4b2, type=rubric, points=2, display={In many samples of 45 adults, an average of 27.9 would support the plan.}]{6.1in}{0.22in}}}}
  \item Calculate $\sigma_X$.\\
        \writeline\gradedhide{\hrulefill}\gradedright{\raisebox{-2pt}{\gradedblank[id=hw4c, type=short, points=1, answer=3.26, check=numeric, tolerance=0.01]{1.3in}{0.2in}}}
  \item Find $P(X \ge 30)$. Show the equality rewrite and the calculator expression.\\
        \writeline\\[1pt]\writeline\gradedhide{\hrulefill}\gradedright{\raisebox{-2pt}{\gradedblank[id=hw4d, type=short, points=2, answer=0.3152, check=numeric, tolerance=0.002]{1.3in}{0.2in}}}
\end{enumerate}

\vspace{8pt}
\textbf{5. Extension:} Use technology or a table to compute $P(X \ge 1)$ for
$X \sim B(20, 0.05)$ two ways: (i) as $1 - P(X = 0)$ and (ii) using
\texttt{binomcdf}. Verify the answers agree and explain why these two approaches
are equivalent.

\vspace{4pt}
\writeline\\[1pt]\writeline\\[1pt]\writeline

\end{practicebox}

\end{document}
